\documentclass[12pt]{amsart}
\usepackage{amsmath,amsthm,amssymb}
\usepackage{dsfont,bm,mathtools}
\usepackage[usenames,dvipsnames,svgnames,table]{xcolor}
\usepackage[top=2.3cm,bottom=2.3cm,left=2cm,right=2cm]{geometry}
\usepackage{hyperref}
\hypersetup{hidelinks}
\allowdisplaybreaks
\emergencystretch=1.5em
\definecolor{ProposalOrange}{RGB}{160,65,0}
\definecolor{RefereePurple}{RGB}{103,34,147}
\newtheorem{thm}{Theorem}[section]
\newtheorem{lem}[thm]{Lemma}
\theoremstyle{remark}
\newtheorem{rem}[thm]{Remark}
\theoremstyle{definition}
\newtheorem{nota}[thm]{Notation}
\definecolor{ResolutionGreen}{RGB}{0,100,55}
\newcommand{\C}{\mathds{C}}
\newcommand{\Q}{\mathds{Q}}
\newcommand{\rv}{\operatorname{rv}}
\newcommand{\Gal}{\operatorname{Gal}}
\newcommand{\Aut}{\operatorname{Aut}}
\newcommand{\RV}{\mathrm{RV}}
\newcommand{\RES}{\mathrm{RES}}
\newcommand{\VF}{\mathrm{VF}}
\newcommand{\K}{\mathds{k}}
\newcommand{\Z}{\mathds{Z}}
\newcommand{\N}{\mathds{N}}
\newcommand{\vrv}{\operatorname{vrv}}
\newcommand{\GL}{\operatorname{GL}}
\newcommand{\dcl}{\operatorname{dcl}}
\newcommand{\pr}{\operatorname{pr}}
\newcommand{\defn}{\mathrm{def}}
\newcommand{\given}{\,\middle|\,}
\newcommand{\Sbase}{\mathbb S}
\newcommand{\RVlang}{\mathcal L_{\mathrm{RV}}}
\numberwithin{equation}{section}
\pagestyle{plain}

\newcommand{\A}{\mathds{A}}
\newcommand{\Gm}{\mathbb G_m}
\newcommand{\Var}{\mathrm{Var}}
\newcommand{\Kres}{K_0(\RES)}
\newcommand{\Bangres}{{!K_0(\RES)}}
\newcommand{\Kraw}{K_0^{\widehat\mu,\mathrm{raw}}(\Var_{\C})}
\newcommand{\Kmon}{K_0^{\widehat\mu}(\Var_{\C})}
\newcommand{\E}{\mathbb E_b}
\newcommand{\Vol}{\operatorname{Vol}}
\newcommand{\vv}{\operatorname{val}}
\newcommand{\LiftingTheorem}{the lifting theorem of Section~3}

\theoremstyle{plain}
\newtheorem*{thm*}{Theorem}

\newcommand{\R}{\mathds R}
\newcommand{\Sfiber}{\mathbb S_b}
\newcommand{\Rbase}[1]{{!K_0(\RES_{#1})}}
\newcommand{\IsVFVG}{(\VF,\Gamma)\text{-generated}}
\newcommand{\gsk}{K_+}
\begin{document}
\setcounter{section}{5}
\setcounter{subsection}{1}
\setcounter{thm}{15}
{\small\noindent Forey--Yin, \emph{Bounded integral and motivic Milnor fiber}.
Section 5 review: R086--R087.
Approved context is black; proposed manuscript text is orange.}
\medskip

\noindent The preceding inflated-family theorem has been incorporated.
Its source pieces and base pieces have separate indices, and its
identity is
\[
 \E\!\left(\int[A^{\mathrm{inf}}]\right)
       =[\Gm]^n\sum_a\chi_b(E_a)v_a .
\]

{\color{ProposalOrange}
\begin{rem}[The polyhedral base pieces]
Each $E_a\subseteq\Q^n$ is the intersection with $\Q^n$ of
a finite Boolean combination of closed rational polyhedra in
$\R^n$. Here a closed rational polyhedron is specified by
finitely many affine inequalities
\[
 a_j\cdot x\geq c_j,\qquad a_j\in\Q^n,\quad c_j\in\Q.
\]
In particular, the $E_a$ need not themselves be closed
polyhedra; $(0,1)\cap\Q$ is a simple example.
The finite arrangement of the bounding hyperplanes refines
them into relatively open rational polyhedral cells, including
cells of lower dimension. Such a refinement preserves the
preceding formula by additivity of $\chi_b$ and constancy of
the coefficients on each original base piece.

The constants $c_j$ may be nonzero. If the chosen defining
conditions are homogeneous, the resulting pieces can instead
be refined into relatively open rational polyhedral cones.
This conditional observation about the value-group formulas
does not change the standing assumption that
$\Gamma(\Sbase)$ is nontrivial.
\end{rem}
}
\medskip
{\normalfont\small\color{RefereePurple}
\noindent\textbf{Referee R086 [Error] -- Polyhedra versus
constructible polyhedral sets.}
A definable value-group set need not be an intersection of closed
half-spaces. Use finite Boolean combinations or relatively open
polyhedral cells, and use inequalities to define half-spaces.
Homogeneous pieces may be relatively open cones.\par}
\medskip

{\color{ProposalOrange}
\subsection{Fubini theorem}
In this subsection, return to an arbitrary
$(\VF,\Gamma)$-generated ground substructure $\Sbase$,
with the standing nontriviality of $\Gamma(\Sbase)$.
All definability remains in the original RV language.
Let $A,A'\in\VF_*$ and let
$f:A\to B$, $f':A'\to B$ be definable maps.
Write $A_b=f^{-1}(b)$, $A'_b=(f')^{-1}(b)$ and
$\Sfiber=\Sbase\langle b\rangle$.
For a parameter base $\mathbb T$, write
\[
 R(\mathbb T)=\Rbase{\mathbb T}
\]
for the ungraded RES quotient computed with
$\mathbb T$-definable sets, maps and value-fiber relations.
The subscript $\mathrm b$ in $\mathbb E_{\mathrm b}$ denotes
the bounded retraction; the parameter base is specified
separately.

\begin{thm}\label{fubini}
Suppose that for every $b\in B$ there are
$\Sfiber$-definable objects
$\bm U_b,\bm U'_b\in\RV[*]$ and
$\Sfiber$-definable bijections
\[
 H_b:A_b\xrightarrow{\ \sim\ }\mathbb L\bm U_b,
 \qquad
 H'_b:A'_b\xrightarrow{\ \sim\ }\mathbb L\bm U'_b
\]
such that
\[
 \mathbb E_{\mathrm b}^{\Sfiber}([\bm U_b])
 =
 \mathbb E_{\mathrm b}^{\Sfiber}([\bm U'_b])
 \quad\text{in }R(\Sfiber).
\]
Then, over the original ground $\Sbase$,
\[
 \mathbb E_{\mathrm b}^{\Sbase}\!\left(\int[A]\right)
 =
 \mathbb E_{\mathrm b}^{\Sbase}\!\left(\int[A']\right)
 \quad\text{in }R(\Sbase).
\]
As usual, $B$ and its points are interpreted in the sufficiently
saturated ambient model.
\end{thm}

The lifting hypotheses mean
\[
 [A_b]_+=[\mathbb L\bm U_b]_+,\qquad
 [A'_b]_+=[\mathbb L\bm U'_b]_+
\]
in the positive Grothendieck semiring over $\Sfiber$.
For these categories of definable sets, finite scissor
decompositions assemble into actual definable bijections;
the displayed maps specify the required witnesses explicitly.
The ground-generation assumption supplies the domain of
the global integral from Section~3.
No generation assumption is imposed on every $\Sfiber$
merely to evaluate the displayed RV retractions.

\begin{rem}\label{rem:volb}
Whenever $\Sfiber$ is $(\VF,\Gamma)$-generated, the fiber
integral of Section~3 is available and the hypothesis can
equivalently be written
\[
 \mathbb E_{\mathrm b}^{\Sfiber}
       \!\left(\int[A_b]\right)
 =
 \mathbb E_{\mathrm b}^{\Sfiber}
       \!\left(\int[A'_b]\right).
\]
This applies, for example, to a base parameter $b$ consisting
of VF and Gamma coordinates.
If $B$ is definably bijective over $\Sbase$ to a set of such
tuples, transport the parameter along that bijection:
the two parameter bases have the same definable closure,
so their definable categories and the stated equalities agree.
This does not assert equality of their generated substructures.
\end{rem}

\emph{Witnesses at the beginning of the proof.}
After the initial partition by the images of $f$ and $f'$,
compactness gives a finite definable partition $(B_i)_i$
such that the two RV representatives, their structure maps
and finite-fiber bounds, and both bijection graphs
$H_b,H'_b$ are all given by fixed formulas on each $B_i$.
Indeed validity of each fixed finite package is a first-order
condition, and these conditions cover the parameter space
in the saturated model. Choose quantifier-free formulas
$\xi_i(x,y)$, $\xi'_i(x,y')$ for the two representative objects.

For the graph families over $B_i$, the replacements are the
explicit definable bijections
\[
 (b,a)\longmapsto(b,H_b(a)),\qquad
 (b,a')\longmapsto(b,H'_b(a')).
\]
The base coordinate and all finite RV labels are retained.
The further preparation and descent of the RES equality
witnesses belong to the subsequent argument.
}
\medskip

{\normalfont\small\color{RefereePurple}
\noindent\textbf{Referee R087 [Clarification; revised] --
Witnesses for the lifting equalities.}
Require positive lifting equalities, equivalently definable
bijections. Otherwise stabilization objects would have to be
included and uniformized. State the ground-generation hypothesis
needed for the global integral, while retaining the correct
fiber coefficient base.\par}
\medskip

{\normalfont\small\color{ProposalOrange}
\noindent\textbf{Proposed disposition of R086--R087.}
Replace the polyhedral remark, clarify the general parameter
scope at the opening of Section~5.2, and state the lifting
bijections directly in Theorem~5.17.
Replace the first paragraph of Remark~5.18 by the displayed
fiber-integral clarification and strengthen the first compactness
paragraph of the proof to retain both bijection graphs.
The arithmetic comparison and unramified-point assertions
(R088--R089), and the later RV-coordinate and witness-descent
arguments, remain for their own review batches.}
\end{document}
